\documentclass[10pt]{article}
\usepackage[margin=20mm]{geometry}
\usepackage{amsmath,amssymb,mathtools,fontspec,unicode-math,xurl,hyperref}
\IfFontExistsTF{DejaVu Serif}{\setmainfont{DejaVu Serif}}{\setmainfont{Latin Modern Roman}}
\setmathfont{latinmodern-math.otf}
\setlength{\parindent}{0pt}\setlength{\parskip}{6pt}\sloppy
\title{Technical Verification of
Explicit Reciprocity Laws and Iwasawa Theory
for Modular Forms}
\author{AI-assisted technical review}
\date{26 October 2026}
\begin{document}
\maketitle
Abstract: Three corrections affect the stated proof: unrestricted uniqueness of a local Kirillov inverse fails, the full cyclotomic Iwasawa algebra at \(p=2\) admits counterexamples to two Fitting-ideal lemmas, and the printed sign projections do not give the claimed interpolation formula. The local inverse needed for the newvector has a verified restricted repair. The analytic projections also have a precise local correction. A complete replacement of the prime-two global argument is not supplied. No counterexample to the intended constant-scaled main theorem is claimed. Scientific adjudication remains pending.\par
\section*{Exact source, theorem, and scope}
The source is Matthew Emerton, Robert Pollack, and Tom Weston, Explicit reciprocity laws and Iwasawa theory for modular forms [paper], arXiv:2210.02013v2, dated 9 October 2024, 52 pages. The assignment associates it with DOI 10.1215/00127094-2024-0073 (https://doi.org/10.1215/00127094-2024-0073). Identity with the journal version of record has not been established. All source locations below are one-based physical PDF pages, agreeing with inspected printed labels.\par
Theorem A concerns a fixed newform of weight \(k\geq2\), critical integers \(0\leq j\leq k-2\), local irreducibility, the no-pole hypothesis at \(j+1\), and the specified Euler-system image hypothesis. It allows one nonzero integral multiplier \(C\), independent of the cyclotomic layer \(n\) and the eligible \(j\). It does not prove the unscaled Mazur–Tate conjecture. The residual irreducibility condition in the global sections has an overbar on \(\rho_f\); it is distinct from local \(p\)-adic irreducibility.\par
A native in-chat helper read all 52 pages, with 17 groups of proof and boundary checks. The lead reviewer read the theorem, decisive definitions, local inverse and pairing arguments, global algebra, normalizations, and final deduction; independently reconstructed the counterexamples and repairs; and inspected their original page images. The source views and computations of the helper are attributed to that helper. No separate external referee or independently verified backend model is claimed.\par
The review treats the cited foundations of \(p\)-adic Hodge theory, local Langlands, local duality, Kato’s Euler system, Greenberg’s structure theorem, and Nekovář–Wingberg duality as external inputs. Their complete original proofs were not independently reconstructed. No author code was executed; absent code is not a mathematical defect. The conclusions concern this exact version, not later corrections or author adoption.\par
\section*{The local inverse needs a specified germ}
Problem 1 (Major: Proposition 6.4.1 is false for arbitrary Kirillov functions).  On pp. 21–22, absence of a pole of \(L(\pi,s)\) at \(s=1\) is said to give a unique smooth \(\xi\) such that its restrictions near zero are Kirillov functions and \[(1-P)\xi=\phi,\qquad (P\xi)(x)=\xi(px).\] The proposition imposes no global support condition on \(\xi\).\par
Counterexample, including the representation hypotheses. Take \(p=3\), and a nontrivial ramified quadratic character \(\eta:\mathbb Q_3^\times\to\mathbb C^\times\) with \(\eta(3)=1\). Put \[\chi_1=|\cdot|_p^{-1}\eta,\qquad
 \chi_2\ \text{unramified},\quad\chi_2(p)=2,\qquad
 \pi=\operatorname{Ind}_{\overline B}^{\mathrm{GL}_2(\mathbb Q_p)}(\chi_1\otimes\chi_2),\] in the source’s unnormalized lower-Borel convention. The ratio \(\chi_1/\chi_2\) is ramified, hence is not an exceptional unramified reducibility ratio. Thus \(\pi\) is irreducible and infinite dimensional. Its only unramified Euler factor, in this convention, is \[L(\pi,s)=(1-2p^{1-s})^{-1};\] the denominator at \(s=1\) is \(-1\).\par
The Fourier realization on p. 19 is \(\phi_f(x)=|x|_p\chi_1(x)\widehat f(x)=\eta(x)\widehat f(x)\). For every integer \(m\), the compactly supported function \(f_m=p^{-m}1_{p^{-m}\mathbb Z_p}\) has Fourier transform \(1_{p^m\mathbb Z_p}\). Consequently \(\eta1_{p^m\mathbb Z_p}\) is a Kirillov function. The same realization includes every compactly supported locally constant function on \(\mathbb Q_p^\times\).\par
Now set \[\phi=\eta1_{\mathbb Z_p^\times},\qquad
 \xi_0=\eta1_{\{\operatorname{ord}_p(x)\leq0\}},\qquad
 \xi_1=\xi_0+\eta.\] Because \(\eta(px)=\eta(x)\), both differences are \[\xi_i(x)-\xi_i(px)
 =\eta(x)\bigl(1_{\{\operatorname{ord}_p(x)\leq0\}}-
                   1_{\{\operatorname{ord}_p(x)+1\leq0\}}\bigr)=\phi(x).\] On \(p^m\mathbb Z_p\), the truncation of \(\xi_0\) is zero if \(m>0\), and is \(\eta(1_{p^m\mathbb Z_p}-1_{p\mathbb Z_p})\) otherwise. The truncation of \(\xi_1\) adds \(\eta1_{p^m\mathbb Z_p}\). They therefore satisfy the required local Kirillov condition. A compact open neighborhood differing from such a ball adds only a compactly supported function on \(\mathbb Q_p^\times\). Yet \(\xi_1-\xi_0=\eta\ne0\). The conditional regular-unit-invariance clause does not apply to this \(\phi\). Moreover, the angular average of \(\eta\) is zero, so even the asserted outer-unit Mellin value does not distinguish the solutions.\par
The proof omits ramified unit types when reducing arbitrary germs to unramified geometric sequences. Its final support argument also invokes a condition on the difference that is not a hypothesis of the proposition.\par
Verified repair for the actual input. This counterexample is not a counterexample to the classes \(c_{n,j}\). For the translated newvector \(d_{n,j}\), sufficiently near zero the additive character is trivial and the germ has regular unit weight \(-j\), times an unramified radial newvector tail. Restrict the chosen solution to that finite-dimensional, shift-stable germ space. The shift eigenvalues are \[r_i=\alpha_i/p^{j+1},\] where the \(\alpha_i\) are the nontrivial local Euler parameters. The no-pole condition at \(j+1\) excludes \(r_i=1\). For a Jordan block \(A=rI+N\), \(N^d=0\), the explicit inverse is \[(I-A)^{-1}=\sum_{h=0}^{d-1}\frac{N^h}{(1-r)^{h+1}}.\] Thus repeated Euler roots cause no obstruction. If the Euler factor is one, the newvector germ is eventually zero.\par
Select this inverse germ and extend to outer annuli by \(\xi(x)=\phi(x)+\xi(px)\). This gives uniqueness within the selected germ normalization; a homogeneous difference has zero germ and therefore vanishes everywhere by the same recursion. Apply the operator to the full \(K_z=j!(tx)^{k-2+\ell-j}\phi_v\), retaining the \(j!\) omitted from the special-case inverse display on p. 25. For the actual \(\ell=2-k\), the weight is \(-j\); the cyclotomic twist used in Section 7 cancels it. This supplies the local inverse required at pp. 28–30, including the newvector case, without asserting the original universal uniqueness. The generalized Kirillov lifting remains the cited Colmez input. The correction affects a central construction, although its intended restricted use is repairable.\par
\section*{The prime-two algebraic boundary}
Problem 2 (Major: Lemmas 9.2.2 and 9.2.3 fail in the stated full algebra).  Section 8.1, p. 37, defines \(\Lambda=\mathcal O_L[[\operatorname{Gal}(\mathbb Q(\mu_{p^\infty})/\mathbb Q)]]\). No odd-prime restriction or integral character-component replacement is stated. Lemma 9.2.2 says that a finitely generated torsion module with no nonzero finite submodule has principal zeroth Fitting ideal. Lemma 9.2.3 says that \(I\subseteq A\), with \(A\) of finite index, implies \(I\subseteq\operatorname{Fitt}_\Lambda(A/I)\). Both assertions have counterexamples at \(p=2\).\par
Let \(\mathcal O/\mathbb Z_2\) be any finite extension, with uniformizer \(\varpi\) and finite residue field \(k\). The actual algebra is \[R=\mathcal O[[X,U]]/(U^2+2U),\qquad
 \mathfrak m=(\varpi,U,X),\] where \(1+U\) represents the order-two element and \(X\) the open procyclic factor \(1+4\mathbb Z_2\). This ring is free of rank two over \(\mathcal O[[X]]\), so \(\varpi\) is a non-zero-divisor.\par
First counterexample. The module \(Y=R/(\varpi,U)=k[[X]]\) is cyclic and is killed by the regular element \(\varpi\), so it is torsion. Its nonzero submodules are nonzero \(k[[X]]\)-ideals and are infinite. Nevertheless, \[\operatorname{Fitt}_R(Y)=(\varpi,U)\] is not principal. Indeed, \(U^2+2U\) belongs to the square of the maximal ideal before quotienting, since \(2\in\varpi\mathcal O\). The classes of \(\varpi,U,X\) are independent in \(\mathfrak m/\mathfrak m^2\). For \(J=(\varpi,U)\), its two generators remain independent in \(J/\mathfrak mJ\), because \(\mathfrak mJ\subseteq\mathfrak m^2\). Nakayama’s lemma gives minimal generator number two. This directly contradicts the explicit principality conclusion, without relying on an ambiguous characteristic-ideal convention.\par
Equation reference: eq:relations. Second counterexample and a complete presentation. Take \(A=\mathfrak m\), \(I=\varpi R\). Then \(I\subset A\) and \(R/A=k\) is finite. The module \[M=A/I=(U,X)\subset
 S=k[[X,U]]/(U^2)\] has generators \([U],[X]\) and relation matrix \[\begin{pmatrix}
 \varpi&0&U+2&-X\\
 0&\varpi&0&U
 \end{pmatrix}.
 \] To check completeness, modulo \(\varpi\) write \(a=a_0+Ua_1\), \(b=b_0+Ub_1\), with coefficients in \(k[[X]]\). The relation \(aU+bX=0\) gives \(b_0X=0\) and \(a_0+b_1X=0\). Thus the syzygies over \(S\) are generated by \((U,0)\) and \((-X,U)\). Lifting adds exactly the two \(\varpi\)-relations in [eq:relations] (\#eq:relations).\par
The maximal minors generate \[\operatorname{Fitt}_R(A/I)
 =(\varpi^2,\varpi U,\varpi(U+2),\varpi X)
 =\varpi\mathfrak m.\] But \(\varpi\notin\varpi\mathfrak m\): canceling the regular element \(\varpi\) would give \(1\in\mathfrak m\). Therefore \(I\not\subseteq\operatorname{Fitt}_R(A/I)\), as opposed to Lemma 9.2.3. Even the stronger no-finite-submodule property holds here: \(M=k[[X]]U\oplus Xk[[X]]\), and multiplication by \(X\) is injective.\par
Impact and limits of repair. Lemma 9.2.3 feeds Lemma 9.2.4; the proof of Theorem 9.1.1 on p. 42 uses both that lemma and the equality of characteristic and Fitting ideals from Lemma 9.2.2. These are false universal inputs in the explicitly allowed algebra, not merely unexplained notation.\par
For odd \(p\), the torsion factor has order \(p-1\), invertible in \(\mathcal O\), and the usual character components are two-dimensional regular power-series rings. In that setting the argument can be justified componentwise. For example, \(A/fR\) has no finite submodule on each nonzero component, agrees with \(R/fR\) at height-one primes, and therefore has Fitting ideal \((f)\). An explicit odd-prime restriction removes these counterexamples.\par
Retaining \(p=2\) requires a replacement argument, such as working over an open procyclic algebra and proving an integral comparison back to the full ring. Additional suitable projective-dimension hypotheses would also require proof for the actual arithmetic modules. The allowed scalar \(C\) may make a replacement possible; this report does not establish a uniform prime-two repair. These examples refute the algebraic lemmas, not Theorem A for an explicit modular form.\par
\section*{Sign projections do not give formula (18)}
Problem 3 (Major: the analytic definitions and interpolation use different normalizations).  On p. 42, writing \(\lambda(a)=\lambda(f,z^j,a,p^n)\), the source sets \[\lambda^\pm(a)=\lambda(a)\pm\lambda(-a),\qquad
 \phi(a)=\lambda^+(a)/\Omega^++\lambda^-(a)/\Omega^-.\] There is neither a factor \(1/2\) nor a factor \((-1)^j\). Formula (18) on p. 43 instead uses the period with sign \((-1)^j\chi(-1)\).\par
Equation reference: eq:literal, eq:mellin. For a primitive \(\chi\), put \(\epsilon=\chi(-1)\) and \(S_\chi=\sum_a\chi(a)^{-1}\lambda(a)\). Reindexing by \(a\mapsto-a\) gives \(\sum_a\chi(a)^{-1}\lambda(-a)=\epsilon S_\chi\). Thus the literal definitions imply \[\chi(\theta_{n,j}^{\rm an})
 =\left(\frac{1+\epsilon}{\Omega^+}
       +\frac{1-\epsilon}{\Omega^-}\right)S_\chi
 =\frac{2S_\chi}{\Omega^\epsilon}.
 \] The Mellin calculation from the displayed integral gives \[S_\chi=
 \frac{\chi(-1)\,j!\,p^{nj}\tau(\chi^{-1})}
      {(-2\pi i)^j}\,L(f,\chi,j+1).
 \] For example, substitute \(z=-a/p^n+iy\), integrate each Fourier term, and use \[\sum_a\chi^{-1}(a)e^{-2\pi i ma/p^n}
 =\chi(-m)\tau(\chi^{-1}).\] Formula (18) is \(S_\chi/\Omega^{(-1)^j\epsilon}\), not [eq:literal] (\#eq:literal). For even \(j\) it misses a factor two; for odd \(j\) the sign label is also exchanged.\par
This is a substantive mismatch, not an assumption that an unspecified central \(L\)-value is nonzero. Nonvanishing examples for the even case are available at the highest critical point of any even weight \(k\geq4\): \(j=k-2\), \(s=k-1>(k+1)/2\). The Ramanujan bound used by the source gives an absolutely convergent, nonzero twisted Euler product there. A primitive character has nonzero Gauss sum, so \(S_\chi\ne0\) in this case.\par
Equation reference: eq:projector. Local correction and propagation. To obtain (18) with the displayed formula for \(\phi\), replace the projections by \[\lambda^\pm_j(a)
 =\frac{\lambda(a)\pm(-1)^j\lambda(-a)}2
 \] and choose the periods consistently with these projections. Then the finite-character calculation gives exactly \(S_\chi/\Omega^{(-1)^j\epsilon}\). The change is to the printed definitions; it must not be silently assumed in reviewing their present form.\par
A fixed factor two alone is harmless for the existence of some multiplier \(C\). That observation does not remove the parity error or validate the comparison on pp. 45–46 as written. The correction must be propagated through the cohomological periods, Kato periods, and comparison constants. The error lies in the central analytic-to-algebraic identification; it does not establish failure of the intended weak main theorem.\par
A separate normalization question remains unproved. The calculation in Proposition 10.2.3 displays ratios \[D_\pm=\langle e,\omega_f\rangle\,
           \Omega_f^\pm/\Omega_{f,K}^\pm,\] while the assertion uses a single scalar \(D\). The permitted choices of the two sign components of the Kato class were not reconstructed sufficiently to disprove that assertion. This is a request for a compatible normalization, classified NeedsEvidence, not a fourth demonstrated error. One may instead formulate a componentwise comparison and clear fixed denominators, but this alone is not the missing full prime-two proof.\par
\section*{Other proof and boundary checks}
Period rings, traces, and the local classes. The Frobenius divisibility argument uses all integer specializations, not just one. The trace compatibility produces the required inverse system. The component convention \(\iota_i^-=\iota_0^-\varphi^{-i}\) gives the support and stabilization directions used in the reciprocity statement. The differential pairing differs from the absolute de Rham pairing by \(p^{-n}\); the dual-exponential comparison contributes \(p^n\). Those factors cancel in Proposition 7.4.1, rather than disappearing without explanation.\par
In Lemma 7.3.1, a unipotent difference lowers the degree in \(e_2\), so \(j+1\) differences kill the degree-\(j\) vector. The factor \(a^{k-2-j}\) combined with \(\det^{2-k}\) gives \(a^{-j}\); the twist by \(j\) removes it. The resulting functional factors through the finite-level dual exponential. Local duality supplies the rational class, and orthogonality puts it in \(H_e\subseteq H_f\). The stronger equality \(H_e=H_f\) invoked under the no-pole hypothesis is a cited input; its every-base-change form was not separately proved.\par
The integral argument works in the image of \(H^1(T(1+j))\) in rational cohomology, not by assuming the integral cohomology is torsion-free. The correct torsion identity uses \(A_1=V(1+j)/T(1+j)\): \[H^1(T(1+j))_{\rm tors}
 =H^0(A_1)/\operatorname{im}H^0(V(1+j)).\] This is the discrete-module identity needed for the annihilator calculation; the printed \(H^0(T)\) expressions on p. 37 should be corrected accordingly. Once the integral pairing is supplied by the chosen local-Langlands lattice, the minimal-\(\varpi\)-valuation argument lowers a hypothetical positive denominator and proves the claimed integrality.\par
Newvectors, recurrences, and exceptional parameters. The local Birch sum must be applied in the twisted model. Its angular factor cancels the twisting character on units. The unmodified-model identity on p. 32 should retain the angular character; its later evaluation at powers of \(p\), with \(\chi(p)=1\), has the printed scalar value. The source explicitly distinguishes a \(p\)-deprived newvector from a primitive newvector when twisting removes ramification. The Dirichlet coefficient \(L\)-series in the global argument must retain that distinction.\par
The corestriction sum at \(n=0\) has \(p-1\) terms and yields \(a-\delta-1\); it is not the degree-\(p\) recurrence. The explicit crystalline calculation separately assumes distinct parameters and nonzero denominators. Coincident parameters in the general construction are handled by the Jordan formula above. The no-pole classification in Lemma 10.1.6 follows by comparing Euler-root absolute values, with the integral critical range \(0\leq j\leq k-2\) retained. Negative or half-integral values of its displayed candidates do not create extra admissible cases.\par
The global algebra after its scope is corrected. For odd \(p\), the reflexive hull is free componentwise. The finite control error \(B_n\) need not vanish: the finite Selmer dual is a quotient of the relevant middle module, which is the direction required by the Fitting-ideal estimate. The finite-level coinvariant argument uses the torsion-free quotient’s lack of invariants. The group-ring coefficient calculation gives the stated involution \(\sigma\mapsto\sigma^{-1}\). For a nontrivial coefficient extension, a compatible \(\mathcal O_L\)-linear dual into \(L/\mathcal O_L\), or a trace/dual-lattice convention, is needed for the printed \(\mathbb Z_p\)-valued pairings. This review does not independently reconstruct all those integral comparison conventions.\par
The final divisibility can be justified without canceling a possible zero divisor in a finite group ring. Write \(F=\operatorname{char}_\Lambda H_P^2\), and use the Euler-system input to write \(\alpha=Fb\) in \(\Lambda\). If \(P_n\) denotes the local pairing sum, then \[\iota(\theta^{\rm alg}_{n,j})=F_nP_n,\qquad
 D\iota(\theta^{\rm an}_{n,j})
   =\alpha_nP_n=b_nF_nP_n
   =b_n\iota(\theta^{\rm alg}_{n,j}).\] Thus no cancellation of \(F_n\) is necessary. For a valid scalar \(D\ne0\), the valuation-ring choice \(C=D\) if \(D\in\mathcal O_L\), and \(C=1\) otherwise, is correct. The underlying Kato divisibility and Selmer-complex duality remain the named external inputs.\par
\section*{Conclusion}
The assessment of the exact version is Major. There are demonstrated failures in universal local and algebraic statements, and an incompatible normalization in the interpolation used by the main comparison. The restricted newvector construction and analytic projections have concrete local repairs. The full \(p=2\) argument still needs replacement; the main constant-scaled theorem has not been disproved.\par
Issue | Assessment | Verified repair | Adoption | Remaining work\par
Local inverse | Major | Selected germ; Jordan inverse; retain \(j!\) | Unknown | State restricted hypothesis; retain cited lifting\par
Prime-two algebra | Major | Odd-prime restriction removes examples | Unknown | Full \(p=2\) replacement unproved\par
Sign projections | Major | Formula [eq:projector] (\#eq:projector) | Unknown | Propagate compatible periods\par
Common scalar | Needs evidence | No complete normalization established | Unknown | Specify sign choices\par
Scientific decision: pending. Human reviews: none. No independent acceptance or author adoption is recorded.\par
9 M. Emerton, R. Pollack, and T. Weston, Explicit reciprocity laws and Iwasawa theory for modular forms, arXiv:2210.02013v2 (9 October 2024), 52 pp. Reviewed PDF SHA256: `00768fa79ba756c779da125aa72c75fe762b6b0cde17987d2c84c6fd91c22967`.\par
\end{document}
